\section{Baker--Akhiezer realizations on a prescribed curve}
\label{sec:general-BA-realization}

We construct solutions for arbitrary positive $L,M,N$ directly from a
smooth curve. The two incident arrays come from ordinary and dual BA
bases normalized at $N$ punctures; their two-point kernel supplies the
$L\times M$ array. Meromorphic interpolation determines the weights
and phase coefficients. Its residue identities give Hamilton's equations
without first reconstructing a matrix spectral problem.
For $N=1$ the basis has a single ordinary/dual pair, giving the elementary
vectors and pair composites of the hyperelliptic specialization.

\subsection{Scalar BA bases and their kernel}
\label{sec:multipoint-BA-data}

Fix a smooth compact connected curve $\mathcal C$ of genus $g$,
distinct punctures $R_k$ with local parameters $s_k$, $1\leq k\leq N$,
and regular ordinary and dual BA divisors $\mathscr D,\mathscr D^\vee$
of degree $g+N-1$. Regularity means that the $N$ normalized leading
values determine a BA section uniquely. At the reference divisor this
is $H^0(\mathcal O(\mathscr D-\sum_{k=1}^{N}R_k))=0$, with the
corresponding dual condition. The normalized bases satisfy
\begin{equation}\label{eq:general-BA-asymptotics}
\begin{aligned}
 \psi_k(P,t)&=e^{-it/s_m}(\delta_{km}+O(s_m)),\\
 \psi_k^\vee(P,t)&=e^{it/s_m}(\delta_{km}+O(s_m))
 \quad(P\to R_m),\quad1\leq k,m\leq N.
\end{aligned}
\end{equation}
These are scalar functions on one curve with shared divisor data:
$k$ selects a normalized section, rather than an independent one-point
solution \cite{Krichever1977,KricheverShiota2011}.
At the reference time both bases sum to one: the constant function
has these leading values, so uniqueness gives
$\sum_{k=1}^{N}\psi_k(P,0)=\sum_{k=1}^{N}\psi_k^\vee(P,0)=1$.
The number $N$ specifies these BA punctures independently of the genus;
they need not be infinities in another plane equation for the curve.
Their fixed duality differential $d\Omega_0$ has zero divisor
$\mathscr D+\mathscr D^\vee$ and leading part $-ds_k/s_k^2$ at each
puncture; simple-pole residues there are allowed.

Choose normalized holomorphic differentials $\boldsymbol\omega$ and put
\begin{equation}\label{eq:multipoint-Abel-direction}
 \boldsymbol\omega=(-\boldsymbol U_k+O(s_k))ds_k\quad(P\to R_k),
 \qquad \boldsymbol U=\sum_{k=1}^{N}\boldsymbol U_k,\qquad
 z(t)=z_0-it\boldsymbol U.
\end{equation}
The normalized second-kind integral $\Xi$ has expansions
$\Xi=s_k^{-1}+\Xi_k+O(s_k)$. Its principal parts give this common time.

For a holomorphic basis $du$ and compatible second-kind basis $dr$,
our half-period convention is $2\omega=\oint_a du$,
$2\omega'=\oint_b du$, $2\eta=-\oint_a dr$,
$2\eta'=-\oint_b dr$. Thus $\mathbb B=\omega^{-1}\omega'$ and
$\varkappa=\eta\omega^{-1}$, with $\varkappa$ symmetric.
Write $\theta=\theta[\delta](\,\cdot\mid\mathbb B)$ for a fixed
characteristic. The theta-exponential sigma convention is
\cite{BEL1997,KorotkinShramchenko2012}
\[
 \sigma(u)=C\exp\!\left(\tfrac12u^{\trans}\varkappa u\right)
 \theta((2\omega)^{-1}u),\qquad C\ne0.
\]
The bases and characteristic fix this convention. Here
$du=\boldsymbol\omega$ and $2\omega=\Id_g$; the explicit examples
state their unnormalized bases. These conventions agree with Generapy
\cite{Generapy}.

Base the Abel map $A$ at $R_1$, and put
$\Delta_k(P)=A(P)-A(R_k)$ and $J_k(P)=\Xi(P)-\Xi_k$.
The theta formulas are
\begin{equation}\label{eq:multipoint-BA-theta}
\begin{aligned}
 \psi_k(P,t)&=a_k(P)e^{-itJ_k(P)}
       \frac{\theta(z(t)+\Delta_k(P))}{\theta(z(t))},\\
 \psi_k^\vee(P,t)&=b_k(P)e^{itJ_k(P)}
       \frac{\theta(z(t)-\Delta_k(P))}{\theta(z(t))},
       \qquad1\leq k\leq N.
\end{aligned}
\end{equation}
The fixed factors are determined by the reference sections:
$a_k=\psi_k(P,0)\theta(z_0)/\theta(z_0+\Delta_k)$ and
$b_k=\psi_k^\vee(P,0)\theta(z_0)/\theta(z_0-\Delta_k)$.
They encode the divisors and puncture normalizations.
Equivalently, put $J_{\sigma,k}=J_k-\boldsymbol U^{\trans}\varkappa\Delta_k$,
$a_{\sigma,k}=a_k e^{-z_0^{\trans}\varkappa\Delta_k-\Delta_k^{\trans}\varkappa\Delta_k/2}$
and
$b_{\sigma,k}=b_k e^{z_0^{\trans}\varkappa\Delta_k-\Delta_k^{\trans}\varkappa\Delta_k/2}$.
The sigma versions are
$\psi_k=a_{\sigma,k}e^{-itJ_{\sigma,k}}\sigma(z+\Delta_k)/\sigma(z)$ and
$\psi_k^\vee=b_{\sigma,k}e^{itJ_{\sigma,k}}\sigma(z-\Delta_k)/\sigma(z)$.
Finitely many additional BA times may be included to obtain an open
Jacobian family; the dot continues to denote this selected first time.

The scalar Cauchy--BA kernel $\Omega(P,Q;t)$ has the fixed ordinary
and dual divisors and diagonal normalization
\begin{equation}\label{eq:general-kernel-normalization}
 \Omega(P,Q)\sim\frac1{\beta_Q(s_Q(P)-s_Q(Q))},\qquad
 d\Omega_0=(\beta_Q+O(s_Q))ds_Q.
\end{equation}
At each puncture,
\begin{equation}\label{eq:general-kernel-infinity}
\begin{aligned}
 \Omega(P,Q)&\sim s_k(P)e^{-it/s_k(P)}\psi_k^\vee(Q)
                                      &&(P\to R_k),\\
 \Omega(P,Q)&\sim-s_k(Q)\psi_k(P)e^{it/s_k(Q)}
                                      &&(Q\to R_k),\quad1\leq k\leq N.
\end{aligned}
\end{equation}
We use the differential-kernel characterization of
\cite[Definition~3.2 and (3.22)]{Bertola2021}, converted to this sign.
Differentiation removes the diagonal principal part, and the leading
values at all punctures give
\begin{equation}\label{eq:general-kernel-rate}
 i\dot\Omega(P,Q)=\sum_{k=1}^{N}\psi_k(P)\psi_k^\vee(Q).
\end{equation}
Subtracting the sum leaves a BA section with zero leading vector, hence
zero. For $N=1$ this is the usual BA-product primitive
\cite{GrinevichOrlov1989,AratynNissimovPacheva1997}.

\subsection{Mode arrays and meromorphic interpolation}
\label{sec:general-BA-interpolation}

Put $K=L+M$. Choose distinct ordinary points $Q_a$, $1\leq a\leq K$,
away from the punctures, fixed poles and zeros of $d\Omega_0$.
For nonzero weights $\kappa_a$ and reciprocal gauges $c_a,e_k$ define
\begin{equation}\label{eq:multipoint-mode-arrays}
\begin{aligned}
 u_{ij}^{\rm BA}&=\frac{\sqrt2\kappa_i\kappa_{L+j}}{c_ic_{L+j}}
                  \Omega(Q_i,Q_{L+j}),&
 \bar u_{ij}^{\rm BA}&=-\sqrt2c_ic_{L+j}\Omega(Q_{L+j},Q_i),\\
 v_{jk}^{\rm BA}&=\frac{c_{L+j}}{e_k}\psi_k(Q_{L+j}),&
 \bar v_{jk}^{\rm BA}&=\frac{\kappa_{L+j}e_k}{c_{L+j}}\psi_k^\vee(Q_{L+j}),\\
 w_{ki}^{\rm BA}&=c_ie_k\psi_k^\vee(Q_i),&
 \bar w_{ki}^{\rm BA}&=\frac{\kappa_i}{c_ie_k}\psi_k(Q_i),
 \quad1\leq i\leq L,\ 1\leq j\leq M,\ 1\leq k\leq N.
\end{aligned}
\end{equation}
The bars are independent partner labels. The kernel rate already gives
the $u$ equations; the other arrays require weighted sums over the
opposite ordinary-point group.

At $Q_a$ choose a vanishing coordinate and write
\begin{equation}\label{eq:general-residue-data}
 \boldsymbol\omega=(E_a+O(s_a))ds_a,\qquad
 d\Omega_0=(\beta_a+O(s_a))ds_a,\qquad r_a=2\kappa_a/\beta_a.
\end{equation}
Seek $h_L,h_M$ with finite simple poles only at their respective
groups, of principal coefficients $r_a$, and expansions
\begin{equation}\label{eq:general-h-normalization}
 h_A=s_k^{-1}+a_{A,k}+O(s_k)\quad(P\to R_k),
 \qquad A\in\{L,M\},\quad1\leq k\leq N.
\end{equation}
One additive constant per function may be fixed. Their local constants
$a_{A,k}$ determine the third-group phase rates.

\begin{lemma}[Meromorphic interpolation]\label{lem:general-BA-interpolation}
The functions exist if and only if
\begin{equation}\label{eq:general-residue-criterion}
 \sum_{i=1}^{L}r_iE_i=\boldsymbol U,\qquad
 \sum_{j=L+1}^{K}r_jE_j=\boldsymbol U.
\end{equation}
They are unique up to additive constants, and
\begin{equation}\label{eq:general-residue-weights}
 \kappa_i=\tfrac12\operatorname{Res}_{Q_i}(h_Ld\Omega_0),\quad1\leq i\leq L,
 \qquad
 \kappa_j=\tfrac12\operatorname{Res}_{Q_j}(h_Md\Omega_0),\quad L+1\leq j\leq K.
\end{equation}
\end{lemma}
\begin{proof}
The finite residues of $h\omega_\alpha$ are $r_a(E_a)_\alpha$;
their sum at the punctures is $-U_\alpha$. Their vanishing is exactly
the compact-curve Mittag--Leffler obstruction \cite{Hartshorne1977}.
Two solutions differ by a global holomorphic function, hence a constant.
\end{proof}
The test is two $g$-dimensional linear systems. Generic groups of
$n\geq g$ points have rank $g$ and $n-g$ free principal coefficients.
Every coefficient must be nonzero to retain its modes. Special smaller
groups obey the same consistency test.

\subsection{The kernel convolution identity}
\label{sec:kernel-verification}

The interpolation functions turn the kernel into weighted BA identities,
coupling the three arrays.

\begin{lemma}[Convolution and weighted BA equations]
\label{lem:general-kernel-convolution}
For $h=h_A$ let $(a_0,a_1)=(1,L)$ or $(L+1,K)$ respectively. Then
\begin{equation}\label{eq:general-kernel-convolution}
 \sum_{k=1}^{N}\psi_k(P)\psi_k^\vee(Q)
 =(h(P)-h(Q))\Omega(P,Q)
 -2\sum_{a=a_0}^{a_1}\kappa_a\Omega(P,Q_a)\Omega(Q_a,Q).
\end{equation}
For $1\leq k\leq N$,
\begin{equation}\label{eq:general-weighted-BA}
\begin{aligned}
 i\dot\psi_k(P)&=(h(P)-a_{A,k})\psi_k(P)
       -2\sum_{a=a_0}^{a_1}\kappa_a\psi_k(Q_a)\Omega(P,Q_a),\\
 i\dot\psi_k^\vee(P)&=-(h(P)-a_{A,k})\psi_k^\vee(P)
       -2\sum_{a=a_0}^{a_1}\kappa_a\psi_k^\vee(Q_a)\Omega(Q_a,P).
\end{aligned}
\end{equation}
\end{lemma}
\begin{proof}
Apply the residue theorem in $T$ to
$h(T)\Omega(P,T)\Omega(T,Q)d\Omega_0(T)$.
The zeros of $d\Omega_0$ cancel its fixed divisor poles. The residues
at $P,Q,Q_a$ are $-h(P)\Omega(P,Q)$, $h(Q)\Omega(P,Q)$ and
$2\kappa_a\Omega(P,Q_a)\Omega(Q_a,Q)$. Each $R_k$ contributes
$\psi_k(P)\psi_k^\vee(Q)$, proving the convolution.
For the first weighted equation the finite principal parts cancel
by \cref{eq:general-residue-data}. At $R_m$ the residual leading
value is $-a_{A,m}\delta_{km}$ before the $a_{A,k}$ correction.
After that correction every leading value vanishes, so BA uniqueness
applies. The dual proof has the reversed diagonal sign.
\end{proof}

\subsection{Realization and finite completion}

For $1\leq i\leq L$, $1\leq j\leq M$, $1\leq k\leq N$,
define the phase coefficients and auxiliary phase rates by
\begin{equation}\label{eq:general-phase-values}
\begin{aligned}
 b_i^L&=-\tfrac12h_M(Q_i),&
 b_j^M&=\tfrac12h_L(Q_{L+j}),&
 b_k^N&=\tfrac12(a_{M,k}-a_{L,k}),\\
 \ell_i^L&=-\tfrac12h_M(Q_i),&
 \ell_j^M&=-\tfrac12h_L(Q_{L+j}),&
 \ell_k^N&=-\tfrac12(a_{M,k}+a_{L,k}).
\end{aligned}
\end{equation}
Use the cycle-preserving transformation
\begin{equation}\label{eq:multipoint-phase-gauge}
 u_{ij}=e^{-it(\ell_i^L-\ell_j^M)}u_{ij}^{\rm BA},\quad
 v_{jk}=e^{-it(\ell_j^M-\ell_k^N)}v_{jk}^{\rm BA},\quad
 w_{ki}=e^{-it(\ell_k^N-\ell_i^L)}w_{ki}^{\rm BA},
\end{equation}
with reciprocal partner factors. Its three rates sum to zero in
every triad, preserving the mixing terms.

\begin{theorem}[Canonical BA realization]\label{thm:general-BA-realization}
For regular multipoint data, \cref{eq:general-residue-criterion} gives
a solution of \cref{eq:canonical-H} through
\cref{eq:multipoint-mode-arrays,eq:general-phase-values,eq:multipoint-phase-gauge}.
Conversely, for $g\geq1$ and $\boldsymbol U\ne0$, these residue
conditions are necessary for the displayed BA ansatz to solve the
canonical equations on an open Jacobian family.
For preassigned phases, compatibility up to cycle-preserving linear
gauges is precisely
\begin{equation}\label{eq:multipoint-phase-compatibility}
 2(b_i^L+b_j^M+b_k^N)
 =h_L(Q_{L+j})-h_M(Q_i)+a_{M,k}-a_{L,k},
 \quad1\leq i\leq L,\ 1\leq j\leq M,\ 1\leq k\leq N.
\end{equation}
\end{theorem}
\begin{proof}
The kernel and weighted identities give the raw equations
\begin{equation}\label{eq:multipoint-raw-rates}
\begin{aligned}
 i\dot u_{ij}^{\rm BA}&=\sqrt2\sum_{k=1}^{N}\bar v_{jk}^{\rm BA}\bar w_{ki}^{\rm BA},\\
 i\dot v_{jk}^{\rm BA}&=(h_L(Q_{L+j})-a_{L,k})v_{jk}^{\rm BA}
              +\sqrt2\sum_{i=1}^{L}\bar u_{ij}^{\rm BA}\bar w_{ki}^{\rm BA},\\
 i\dot w_{ki}^{\rm BA}&=(-h_M(Q_i)+a_{M,k})w_{ki}^{\rm BA}
              +\sqrt2\sum_{j=1}^{M}\bar u_{ij}^{\rm BA}\bar v_{jk}^{\rm BA},
\end{aligned}
\end{equation}
and negative partner equations. The phase transformation gives
\cref{eq:canonical-rates}. Comparing sums of the linear rates around
each triad proves \cref{eq:multipoint-phase-compatibility}; admissible
redistributions are the constant gauges of \cref{sec:canonical-phase-gauges}.

For necessity let $\mathsf H_\theta$ be the Hessian of $\log\theta$
and $\partial_v=\sum_{\alpha=1}^{g}v_\alpha\partial_{z_\alpha}$
for a constant direction $v$.
Fay's two-kernel identity is
\begin{equation}\label{eq:general-Fay-power}
 \Omega(P,Q)\Omega(Q,P)d\Omega_0(P)d\Omega_0(Q)
 =-B(P,Q)-\boldsymbol\omega(P)^{\trans}\mathsf H_\theta\boldsymbol\omega(Q),
\end{equation}
where $B$ is the normalized fundamental bidifferential
\cite[p.~26]{Fay1973}; see also \cite[(3.12)]{Korotkin2003}.
Its diagonal time derivative gives
\begin{equation}\label{eq:general-elementary-power}
 \left(\sum_{k=1}^{N}\psi_k(P)\psi_k^\vee(P)\right)d\Omega_0(P)
 =d\Omega_1(P)+\boldsymbol\omega(P)^{\trans}\mathsf H_\theta\boldsymbol U,
\end{equation}
where $d\Omega_1=d\Xi$ has the chosen puncture principal parts.
For $V_M=\sum_{j=L+1}^{K}r_jE_j-\boldsymbol U$, the first-group
charge divided by $\kappa_i$ is
$\sum_{k=1}^{N}\psi_k(Q_i)\psi_k^\vee(Q_i)
+2\sum_{j=L+1}^{K}\kappa_j\Omega(Q_i,Q_j)\Omega(Q_j,Q_i)$.
Its product with $d\Omega_0(Q_i)$ is a fixed coefficient minus
$\boldsymbol\omega(Q_i)^{\trans}\mathsf H_\theta V_M$.
Conservation implies
$\partial_{\boldsymbol U}\partial_{E_i}\partial_{V_M}\log\theta=0$.
The leading cubic pole is the product of the three normal derivatives
of $\theta$. Its irreducible reduced ample divisor has finite
translation stabilizer, so no nonzero constant direction is everywhere
tangent. Holomorphic differentials have no common zero, so $E_i\ne0$;
the nonzero $\boldsymbol U$ hypothesis then forces $V_M=0$.
A column charge gives the other residue equation. The weighted equations
then fix the phase compatibility.
\end{proof}
Sufficiency applies also to genus zero and stationary Abel directions;
the nonzero-direction hypothesis is used only for necessity.

\begin{corollary}[Including prescribed evaluations]\label{cor:BA-finite-list}
Fix regular multipoint data and any finite list of ordinary evaluation
points in the two groups with nonzero principal coefficients.
Adjoining $g+1$ generic points to each group permits both residue
equations to be solved with all new coefficients nonzero.
If all $N$ regular constants of each interpolation function are also
prescribed, $g+N$ generic additional points per group suffice.
The enlarged arrays have the realization of
\cref{thm:general-BA-realization}.
\end{corollary}
\begin{proof}
For the residue equations choose $g+1$ columns such that every $g$
are independent. The solution set for any residual right side is an
affine line whose kernel direction has no zero coordinate. Avoid its
finitely many zero-coordinate parameter values.
For prescribed regular constants the obstruction space is
$H^1(\mathcal O(-\sum_{k=1}^{N}R_k))$, of dimension $g+N-1$.
Its dual consists of holomorphic and third-kind differentials with
at most simple puncture poles. Generic $g+N$ evaluation columns have
the same full-rank and nonzero-kernel property. Apply the same argument
to each group. In genus zero the residue system is empty.
\end{proof}

\section{The one-point specialization}
\label{sec:BA-one-point-specialization}

The general construction becomes a vector--matrix--vector system
when $N=1$. We give its elementary-mode notation, specialize the
higher-flow phases and explain its constrained-KP interpretation.
The curve remains arbitrary throughout this section.

\subsection{Elementary modes and their canonical equations}
\label{sec:BA-one-point}

For $N=1$ put $R_1=P_\infty$, $s_1=s$ and $e_1=1$.
The divisors have degree $g$ and there is a single ordinary/dual pair
$\psi,\psi^\vee$. Choose reference normalization
$\psi(P,0)=\psi^\vee(P,0)=1$, and subtract the constant term of
each $h$ at $P_\infty$. Thus $a_{L,1}=a_{M,1}=0$.
This member has two elementary vectors and a pair-mode matrix.

The combined index and number of canonical pairs are
\begin{equation}\label{eq:BA-mode-counts}
 K=L+M,\qquad p=K+LM.
\end{equation}
We use single indices for the $K$ elementary pairs and double indices
for the $LM$ cross-group pair composites. Their dictionary with the
$H_0$ modes of \cref{eq:canonical-phase-gauge} is
\begin{equation}\label{eq:BA-mode-dictionary}
 \xi_i=\tilde w_{1i},\qquad
 \xi_{L+m}=\tilde v_{m1},\qquad
 \xi_{i,L+m}=\tilde u_{im},
 \quad1\leq i\leq L,\quad1\leq m\leq M,
\end{equation}
with the same rule for barred partners. Put
\begin{equation}\label{eq:BA-phase-dictionary}
 b_i=2b_i^L+b_1^N,\qquad b_{L+m}=2b_m^M+b_1^N,
 \quad1\leq i\leq L,\quad1\leq m\leq M.
\end{equation}
The inherited brackets satisfy $\{\xi_k,\bar\xi_r\}=-i\delta_{kr}$, with the corresponding rule for pair modes, and
\begin{equation}\label{eq:BA-H0-flat}
 H_0=\sum_{k=1}^{K}b_k\xi_k\bar{\xi}_k
 +\sqrt2\sum_{i=1}^{L}\sum_{j=L+1}^{K}
 (\xi_i\xi_j\xi_{ij}+\bar{\xi}_i\bar{\xi}_j\bar{\xi}_{ij}).
\end{equation}


The equivalent one-point theta and sigma formulas, with
$\Xi=s^{-1}+O(s)$ and
$\Xi_\sigma=\Xi-\boldsymbol U^{\trans}\varkappa A(P)$, are
\begin{equation}\label{eq:general-BA-theta}
\begin{aligned}
 \psi(P,t)
 &=e^{-it\Xi(P)}
 \tfrac{\theta(z(t)+A(P))\theta(z_0)}
 {\theta(z(t))\theta(z_0+A(P))}
 =e^{-it\Xi_\sigma(P)}
 \tfrac{\sigma(z(t)+A(P))\sigma(z_0)}
 {\sigma(z(t))\sigma(z_0+A(P))},\\[4pt]
 \psi^\vee(P,t)
 &=e^{it\Xi(P)}
 \tfrac{\theta(z(t)-A(P))\theta(z_0)}
 {\theta(z(t))\theta(z_0-A(P))}
 =e^{it\Xi_\sigma(P)}
 \tfrac{\sigma(z(t)-A(P))\sigma(z_0)}
 {\sigma(z(t))\sigma(z_0-A(P))}.
\end{aligned}
\end{equation}
In this notation the unphased BA modes are
\begin{equation}\label{eq:general-elementary-modes}
\begin{aligned}
 \xi_i&=c_i \psi^\vee(Q_i),&\bar{\xi}_i&=\kappa_i \psi(Q_i)/c_i,
 &&1\leq i\leq L,\\
 \xi_j&=c_j \psi(Q_j),&\bar{\xi}_j&=\kappa_j \psi^\vee(Q_j)/c_j,
 &&L+1\leq j\leq K,
\end{aligned}
\end{equation}
and
\begin{equation}\label{eq:general-pair-modes}
 \xi_{ij}=\frac{\sqrt2\kappa_i\kappa_j}{c_ic_j}\Omega(Q_i,Q_j),
 \qquad \bar{\xi}_{ij}=-\sqrt2c_ic_j\Omega(Q_j,Q_i),
 \quad1\leq i\leq L<j\leq K.
\end{equation}

An \emph{elementary mode} has the shift of one curve-point lift
relative to $P_\infty$. A \emph{pair composite} has the difference
of the two lifts; in the hyperelliptic sheet convention this becomes
the sum of the elementary shifts. The terms specify the shift data,
not a factorization at a common $z$. For $N>1$, the basis sections
instead have shifts $A(P)-A(R_k)$; all sections share the same curve
and kernel, but generally no single base point makes both arrays
one-point elementary vectors.

The raw phase values and equations reduce to
\begin{equation}\label{eq:one-point-phase-values}
 b_i=-h_M(Q_i),\quad1\leq i\leq L,\qquad
 b_j=h_L(Q_j),\quad L+1\leq j\leq K,
\end{equation}
\begin{equation}\label{eq:general-network-rates}
\begin{aligned}
 i\dot\xi_i&=b_i\xi_i+\sqrt2\sum_{j=L+1}^{K}\bar{\xi}_j\bar{\xi}_{ij},&
 i\dot{\bar{\xi}}_i&=-b_i\bar{\xi}_i-\sqrt2\sum_{j=L+1}^{K}\xi_j\xi_{ij},
 &&1\leq i\leq L,\\
 i\dot\xi_j&=b_j\xi_j+\sqrt2\sum_{i=1}^{L}\bar{\xi}_i\bar{\xi}_{ij},&
 i\dot{\bar{\xi}}_j&=-b_j\bar{\xi}_j-\sqrt2\sum_{i=1}^{L}\xi_i\xi_{ij},
 &&L+1\leq j\leq K,\\
 i\dot\xi_{ij}&=\sqrt2\bar{\xi}_i\bar{\xi}_j,&
 i\dot{\bar{\xi}}_{ij}&=-\sqrt2\xi_i\xi_j,
 &&1\leq i\leq L<j\leq K.
\end{aligned}
\end{equation}
These are Hamilton's equations for \cref{eq:BA-H0-flat}.
For this phase gauge, write $\mathsf A_L=\diag(b_1,\ldots,b_L)$,
$\mathsf A_M=\diag(b_{L+1},\ldots,b_K)$.
An optional scalar spectral origin $\rho$ gives
\begin{equation}\label{eq:BA-H0-pencil}
\begin{aligned}
 \mathsf L_0(\nu)&=\begin{pmatrix}
 \rho\Id_L-\mathsf A_L+\nu\Id_L&\sqrt2\boldsymbol\xi_{LM}&\bar{\boldsymbol\xi}_L\\
 \sqrt2\bar{\boldsymbol\xi}_{LM}^{\trans}&\rho\Id_M-\mathsf A_M-\nu\Id_M&\boldsymbol\xi_M\\
 -\boldsymbol\xi_L^{\trans}&-\bar{\boldsymbol\xi}_M^{\trans}&\rho
 \end{pmatrix},\\
 \mathsf T_0(\nu)&=\begin{pmatrix}
 0&0&\bar{\boldsymbol\xi}_L\\0&0&-\boldsymbol\xi_M\\
 -\boldsymbol\xi_L^{\trans}&\bar{\boldsymbol\xi}_M^{\trans}&-\nu
 \end{pmatrix}.
\end{aligned}
\end{equation}
Here $\boldsymbol\xi_L,\boldsymbol\xi_M$ are the column vectors of
the two elementary groups and $\boldsymbol\xi_{LM}$ is their pair
matrix. This is the one-point form of the canonical pencil after
its phase gauge; $\rho$ only shifts the spectral ordinate.

\subsection{Higher-flow phases in the one-point notation}
\label{sec:BA-one-point-higher-flows}

We specialize \cref{thm:BA-higher-flows} using the modes and pencil
just defined. This gives a modal-power expression for the finite
phase constants of a higher Hamiltonian.

For $N=1$, choose the ordinary gap complement at $P_\infty$.
Then $p_1=\sum_{\alpha=1}^{g}v_\alpha s^{-n_\alpha}$ for its
Weierstrass gaps $n_\alpha$, and $\varphi_1^N=0$.
The elementary phase rates are
$b_i^{(r,\nu_0)}=-\varphi_i^L$,
$b_{L+j}^{(r,\nu_0)}=\varphi_j^M$; pair phases are the opposite
sums of these rates. For example, in the one-point $H_0$ pencil
\cref{eq:BA-H0-pencil}, with optional spectral origin $\rho$, the
fourth trace has $q=\lambda^3/\nu$. Put $\epsilon_a=\rho-b_a$,
$\chi_a=1$ in the first ordinary group and $\chi_a=-1$ in the second.
Write $I_a=\xi_a\bar\xi_a$, $1\leq a\leq K$, and
$I_{ij}=\xi_{ij}\bar\xi_{ij}$, $1\leq i\leq L<j\leq K$.
For the conserved power combinations
$C_i=I_i-\sum_{j=L+1}^{K}I_{ij}$ and
$C_j=I_j-\sum_{i=1}^{L}I_{ij}$, perturbation at infinity gives
\[
 \lambda=\chi_a\nu+\epsilon_a-\chi_aC_a/\nu+O(\nu^{-2}),\qquad
 \operatorname{reg}_{Q_a}q=3\chi_a(\epsilon_a^2-C_a).
\]
The phase is this constant plus the finite part of $\mathfrak f$,
with the elementary sign just specified. Thus the one-point
modal-power interpretation is retained inside the multipoint theorem.

\subsubsection{Scalar constrained operators and conserved coefficients}

For $N=1$, the additive constants of $h_L,h_M$ may be chosen so that
$\mathsf C_L=\mathsf C_M=0$. The matrix operators in
\cref{eq:BA-constrained-operators} are then scalar, with
$\boldsymbol\psi=\psi$ and $\boldsymbol\psi^\vee=\psi^\vee$.
This specialization also exposes two conserved quantities.
Write $I_k=\xi_k\bar\xi_k$, $1\leq k\leq K$. Expanding $\partial^{-1}f=f\partial^{-1}-(\partial f)\partial^{-2}
+O(\partial^{-3})$ and using \cref{eq:general-network-rates} gives
\begin{equation}\label{eq:BA-constrained-H-coefficient}
 \mathcal A_L-\mathcal A_M
 =2\left(\sum_{\ell=1}^{L}I_\ell-\sum_{j=L+1}^{K}I_j\right)\partial^{-1}
      -2H_0\partial^{-2}+O(\partial^{-3}).
\end{equation}
The first two negative-order coefficients of the commutator are
the derivatives of these coefficient differences, up to sign.
Their vanishing recovers conservation of the group power difference
and $H_0$ on the BA realization.

\subsection{The prescribed spectral curve and BA reconstruction}
\label{sec:BA-spectral-converse}

The interpolation functions $h_L,h_M$ give spectral coordinates on
the prescribed curve. With several punctures, identifying that curve
with the normalized characteristic curve requires the birationality
condition stated below.

\begin{proposition}[The spectral curve of a prescribed realization]
\label{prop:BA-prescribed-spectral-curve}
For the phase coefficients \cref{eq:general-phase-values}, put
\begin{equation}\label{eq:BA-prescribed-spectral-map}
 \nu(P)=-\tfrac12(h_L(P)+h_M(P)),\qquad
 \lambda(P)=\tfrac12(h_M(P)-h_L(P)).
\end{equation}
The characteristic polynomial of the canonical pencil vanishes on
this map. If the map is birational onto its image, $\mathcal C$
normalizes the full characteristic curve. A sufficient condition is
that one puncture has a finite limiting $\lambda$ distinct from the
other punctures, in particular pairwise distinct $b_k^N$.
For $N=1$ this condition is automatic.

Let $\Gamma$ be the projective plane characteristic curve. Write
$d=L+M+N$ and $G=LM+MN+NL-d+1$. On this birational locus
the total singularity defect beyond the ordinary three infinity
multiple points is
\begin{equation}\label{eq:BA-prescribed-genus-defect}
 \sum_{q\in\operatorname{Sing}(\Gamma)}\delta_q
 -\binom L2-\binom M2-\binom N2=G-g.
\end{equation}
For distinct phases within groups the additional defect is affine.
\end{proposition}
\begin{proof}
Let $\mathsf X_{\rm BA}$ be the block matrix in
\cref{eq:canonical-pencil} with the raw arrays
\cref{eq:multipoint-mode-arrays} and the phase diagonals
\cref{eq:general-phase-values}. The explicit vector is
\begin{equation}\label{eq:BA-prescribed-eigenvector}
 r(P)=\begin{pmatrix}
 (-\kappa_i\Omega(Q_i,P)/c_i)_{i=1}^{L}\\
 (c_{L+j}\Omega(Q_{L+j},P))_{j=1}^{M}\\
 (e_k\psi_k^\vee(P))_{k=1}^{N}
 \end{pmatrix}.
\end{equation}
The first two block rows of
$(\mathsf X_{\rm BA}+\nu(P)\mathsf D)r=\lambda(P)r$
are the convolution identities for $h_M,h_L$. The third follows by
subtracting their dual weighted equations:
\[
 \sum_{i=1}^{L}\kappa_i\psi_k^\vee(Q_i)\Omega(Q_i,P)
 -\sum_{j=L+1}^{K}\kappa_j\psi_k^\vee(Q_j)\Omega(Q_j,P)
 =(\lambda(P)-b_k^N)\psi_k^\vee(P).
\]
The phase transformation \cref{eq:multipoint-phase-gauge} conjugates
the pencil and vector by the same diagonal matrix, preserving this
spectral equation. Adding a scalar origin $\rho\Id_d$ to the
pencil shifts $\lambda$ by $\rho$ and preserves the eigenvectors.

The poles of $\nu$ are simple, at all $Q_a,R_k$. Their branches obey
\[
 \lambda-\nu\to-2b_i^L\ (P\to Q_i),\qquad
 \lambda+\nu\to-2b_j^M\ (P\to Q_{L+j}),\qquad
 \lambda\to b_k^N\ (P\to R_k).
\]
If one $R_k$ has a unique finite limiting ordinate, it is the sole
point over its image in the normalization. Its simple pole of $\nu$
makes this point unramified, so the covering degree is one.
The same argument applies to a uniquely separated outer branch.
For a birational map, $\deg\nu=d$ implies that the minimal polynomial
of $\lambda$ over $\C(\nu)$ has degree $d$. The monic characteristic
polynomial has that degree and vanishes, so it is this minimal
polynomial. The plane degree is $d$; subtracting the ordinary
infinity defects gives \cref{eq:BA-prescribed-genus-defect}.
\end{proof}

\paragraph{Prescribed curves and spectral fibres.}
Fix the marked spectral curve and its functions $\lambda,\nu$,
with distinct phase intercepts. On the regular framed spectral
correspondence, a smooth genus-$G$ level has $G$ divisor-class
parameters and $d-1$ independent mode gauges, hence dimension
$G+d-1=p$. A prescribed genus-$g$ normalization instead supplies
$g+d-1$ parameters after identifying equivalent meromorphic frames.
Relative to an expected-dimensional spectral fibre, this realized
family has codimension $G-g$. It is an invariant normalization-based
family on a special level when $g<G$, and is open in the regular
fibre when $g=G$.

This family has a concrete spectral-sheaf interpretation. Separate
the forced ordinary infinity points, and let
$n:\mathcal C\to\Gamma_{\mathrm{sp}}$ normalize the remaining
integral spectral model. At the reference time the components of
\cref{eq:BA-prescribed-eigenvector} are sections of
\[
 \mathscr L=\mathcal O\left(\mathscr D^\vee+
                         \sum_{a=1}^{K}Q_a\right).
\]
Subtracting the full pole fibre of $\nu$ gives
$\mathscr L(-\nu^{-1}(\infty))=\mathcal O(\mathscr D^\vee-
\sum_{k=1}^{N}R_k)$, with both cohomology groups zero.
Its leading-value sequence shows that the $d$ components form a
basis and that $\nu_*\mathscr L\simeq\mathcal O^d$.
Multiplication by $\lambda$ gives the pencil, so the component
spectral sheaf on $\Gamma_{\mathrm{sp}}$ is $n_*\mathscr L$.
Thus the realized family is the stratum obtained by pushing line
bundles forward from the normalization, with the mode framings
retained. Its divisor-class dimension is $g$.
The full compactified spectral fibre can also contain sheaves with
singularity gluing data \cite{AdamsHarnadHurtubise1990II,Beauville1990}.
The general solution of
\cref{thm:BA-spectral-converse} reconstructs the curve belonging to
generic initial modes. For repeated phases its active dimensions,
moment data and restored block bases are accounted for separately
in Appendix~\ref{app:BA-active-converse}.

With several punctures, the interpolation functions can instead factor
through a covering; curve identification then requires birationality.
The stronger completion in \cref{cor:BA-finite-list} permits distinct
differences $a_{M,k}-a_{L,k}$ to be prescribed, giving the sufficient
branch separation used in the proposition.
For $N=1$, the product coordinates $\lambda-\nu,\lambda+\nu$ have
degrees $M+1,L+1$; their image has bidegree $(L+1,M+1)$ and arithmetic
genus $LM$, recovering the one-point defect $LM-g$.

\begin{theorem}[Generic BA reconstruction]\label{thm:BA-spectral-converse}
Fix positive $L,M,N$ and arbitrary phase coefficients. Generic
canonical initial modes reconstruct scalar multipoint BA bases and
their kernel on a smooth normalized active spectral curve.
For distinct phases within groups all directions are active and the
genus is $G=LM+MN+NL-L-M-N+1$; the last group supplies $N$
punctures, and the ordinary and dual divisors have degree $G+N-1$.
For repeated phases, use constant canonical bases within equal-phase
blocks and remove inactive directions as in
Appendix~\ref{app:BA-active-converse}. Apply the construction with
active dimensions $L_a,M_a,N_a$ and genus
\cref{eq:BA-active-genus}, then restore the zero directions and bases.
The original coordinates can then be constant linear combinations
of the active scalar BA evaluations.
\end{theorem}
\begin{proof}
Work in the $H_0$ phase gauge of \cref{eq:canonical-phase-gauge},
omitting its tildes during this proof. The pencil retains the diagonal
blocks $-2\mathsf B_L,-2\mathsf B_M,\mathsf B_N$, and its time matrix is
\begin{equation}\label{eq:multipoint-H0-time}
 \mathsf T_0(\nu)=
 \begin{pmatrix}0&0&\bar{\boldsymbol w}^{\trans}\\
 0&0&-\boldsymbol v\\
 -\boldsymbol w&\bar{\boldsymbol v}^{\trans}&-\nu\Id_N
 \end{pmatrix}.
\end{equation}
For distinct phases the regular normalized curve has branches
\begin{equation}\label{eq:BA-converse-infinity}
\begin{aligned}
 \lambda&=\nu-2b_i^L+O(\nu^{-1})&&\text{at }Q_i,\quad1\leq i\leq L,\\
 \lambda&=-\nu-2b_j^M+O(\nu^{-1})&&\text{at }Q_{L+j},\quad1\leq j\leq M,\\
 \lambda&=b_k^N+O(\nu^{-1})&&\text{at }R_k,\quad1\leq k\leq N.
\end{aligned}
\end{equation}
The nonempty regular locus, including repeated-phase active
normalizations, is proved in \cref{prop:BA-active-regularity}.
Put $s_k^{-1}=-\nu$ at $R_k$. The functions
\begin{equation}\label{eq:BA-converse-h}
 h_L=-\lambda-\nu,\qquad h_M=\lambda-\nu
\end{equation}
have the required poles and satisfy
$h_M(Q_i)=-2b_i^L$, $h_L(Q_{L+j})=2b_j^M$,
$a_{L,k}=-b_k^N$, $a_{M,k}=b_k^N$.
The interpolation and phase conditions thus come from the initial
spectral curve, rather than restricting initial modes further.

Let $f_k$ be the last $N$ coordinate vectors, $\mathsf J$ their
coordinate projection, and $a=\sum_{k=1}^{N}f_k$.
Choose right and left spectral eigenvectors $r,\ell$ with
$i\dot r=-\mathsf T_0r$, $i\dot\ell=\ell\mathsf T_0$ and reference
normalizations $a^{\trans}r(P,0)=\ell(P,0)a=1$. Define
\begin{equation}\label{eq:BA-converse-functions}
 \psi_k(P)=\ell(P)f_k,\qquad
 \psi_k^\vee(P)=f_k^{\trans}r(P),\qquad
 \Omega(P,Q)=-\frac{\ell(P)r(Q)}{\nu(P)-\nu(Q)}.
\end{equation}
Since $\mathsf T_0(\nu(P))-\mathsf T_0(\nu(Q))
=-(\nu(P)-\nu(Q))\mathsf J$, this kernel satisfies
\cref{eq:general-kernel-rate}. The eigenvalue equations remove apparent
poles at distinct points of an equal-$\nu$ fiber. At $R_k$, the leading
eigenvectors are its coordinate direction and the temporal equation
gives exactly \cref{eq:general-BA-asymptotics}.

For the divisor count use the positive-degree spectral component
line bundle $\mathscr L$, whose pushforward under $\nu$ is
$\mathcal O^d$ \cite{Beauville1990}. Its degree is $G+d-1$.
The last $N$ components vanish simply at the $L+M$ outer infinity
points. Their remaining common BA divisor has degree $G+N-1$.
More explicitly, put $\mathscr E=\mathscr L(-\nu^{-1}(\infty))$.
Then $H^0(\mathscr E)=H^1(\mathscr E)=0$, since its pushforward is
$\mathcal O(-1)^d$. Adding the $N$ punctures gives an exact
leading-value isomorphism
$H^0(\mathscr E(\sum_{k=1}^{N}R_k))\simeq\C^N$.
This proves the required BA uniqueness. The transpose module gives
the dual statement; the analytic infinity justification for repeated
phases is given in Appendix~\ref{app:BA-active-converse}.
The duality differential is
\begin{equation}\label{eq:BA-converse-duality}
 d\Omega_0=-\frac{d\nu}{\ell(P)r(P)}.
\end{equation}
It is regular at the outer branches, has the required double poles
at $R_k$, and its apparent projection-branch poles cancel with $d\nu$.

The reference normalization gives a concrete recovery chart. Set
\begin{equation}\label{eq:BA-converse-constants}
\begin{aligned}
 c_i&=\sum_{k=1}^{N}w_{ki}(0),&
 \bar c_i&=\sum_{k=1}^{N}\bar w_{ki}(0),&
 \kappa_i&=c_i\bar c_i,\quad1\leq i\leq L,\\
 c_{L+j}&=\sum_{k=1}^{N}v_{jk}(0),&
 \bar c_{L+j}&=\sum_{k=1}^{N}\bar v_{jk}(0),&
 \kappa_{L+j}&=c_{L+j}\bar c_{L+j},\quad1\leq j\leq M.
\end{aligned}
\end{equation}
The barred constants are independent partner normalizations, not
complex conjugates; the weights are products of initial sums.
On the open chart where these sums are nonzero, outer eigenvector
expansions give, for $1\leq i\leq L$, $1\leq j\leq M$ and $1\leq k\leq N$,
\begin{equation}\label{eq:BA-converse-recovery}
\begin{aligned}
 \psi_k^\vee(Q_i)&=w_{ki}/c_i,&
 \psi_k(Q_i)&=\bar w_{ki}/\bar c_i,\\
 \psi_k(Q_{L+j})&=v_{jk}/c_{L+j},&
 \psi_k^\vee(Q_{L+j})&=\bar v_{jk}/\bar c_{L+j},\\
 \Omega(Q_i,Q_{L+j})&=\frac{u_{ij}}{\sqrt2\bar c_i\bar c_{L+j}},&
 \Omega(Q_{L+j},Q_i)&=-\frac{\bar u_{ij}}{\sqrt2c_ic_{L+j}}.
\end{aligned}
\end{equation}
For example $r_i\sim-\nu/c_i$ at $Q_i$, and the last eigenvector
rows give $f_k^{\trans}r=w_{ki}/c_i$. The opposite outer block then
gives the kernel recovery, with slope difference two. At an outer
point, in coordinate $1/\nu$, the relevant $h$ has principal
coefficient $-2$ and $d\Omega_0$ has coefficient $-\kappa_a$, so
the recovered weights agree with \cref{eq:general-residue-weights}.
Thus \cref{eq:multipoint-mode-arrays} recovers every mode, with
$e_k=1$. Weighted reference anchors give other generic charts.
The forward theorem and uniqueness of the polynomial ODE identify
the evolution. Restore \cref{eq:canonical-phase-gauge} for $H$.

For repeated phases apply this argument in the constant active bases
of Appendix~\ref{app:BA-active-converse}. Their resolved infinity
directions supply the same normalized leading values. Restore
inactive zero directions and invert the canonical bases afterwards.
\end{proof}

For comparison with a prescribed realization, its fixed $e_k$ gauges
are set to one by the puncture-vertex canonical gauge transformations.
The common rescaling of \cref{sec:general-BA-interpolation} accounts
for the single redundancy among all vertex normalization constants.
The reference partition of unity in \cref{sec:multipoint-BA-data}
then makes \cref{eq:BA-converse-constants} recover the original
weights and gauges directly.

\subsection{Higher Hamiltonians and one-point BA times}
\label{sec:BA-higher-flows}

The commuting Hamiltonians of \cref{sec:spectral-flows} act on
the same BA data in different directions. We now compute those
directions and the accompanying constant mode phases. Jacobian
linearization and its description by principal parts are classical
\cite[\S\S6--7]{Griffiths1985}; see also
\cite{Beauville1990,AdamsHarnadHurtubise1993}.
Here a gap normalization converts the direction into a one-point BA
time, and finite parts give the explicit phases of the canonical
modes. The construction applies to every prescribed realization,
including coincident phases and singular plane spectral images.

Fix a complex spectral value $\nu_0$ and an integer $r\geq2$. On the
$H_0$ pencil \cref{eq:BA-H0-pencil}, put
\begin{equation}\label{eq:BA-shifted-flow}
\begin{aligned}
 H_{r,\nu_0}&=\tfrac1r\tr(\mathsf X+\nu_0\mathsf D)^r,
 &\mathsf G&=(\mathsf X+\nu_0\mathsf D)^{r-1},\\
 \mathsf T_{r,\nu_0}(\nu)
 &=\frac{\mathsf G-\mathsf L_0(\nu)^{r-1}}{\nu-\nu_0},
 &q(P)&=\frac{\lambda(P)^{r-1}}{\nu(P)-\nu_0}.
\end{aligned}
\end{equation}
The quotient defining $\mathsf T_{r,\nu_0}$ is polynomial, and
$i\dot{\mathsf L}_0=[\mathsf L_0,\mathsf T_{r,\nu_0}]$.
Indeed $[\mathsf L_0,\mathsf G]=(\nu-\nu_0)[\mathsf D,\mathsf G]$.
All coefficients of $H_{r,\nu_0}$ in $\nu_0$ are commuting spectral
Hamiltonians. The unshifted traces are $H_{r,0}$.

For $g\geq1$, let $n_1<\cdots<n_g\leq2g-1$ be the Weierstrass
gaps at $P_\infty$: these are the positive integers that cannot be
the pole order of a function having its only pole there. Write
$\operatorname{pp}_Q$ for the principal part at a point $Q$.

\begin{lemma}[One-point normalization]\label{lem:BA-higher-normalization}
There is a unique meromorphic function $\mathfrak f$ with poles only
at $Q_1,\ldots,Q_K,P_\infty$ such that
\[
 \operatorname{pp}_{Q_k}\mathfrak f=-\operatorname{pp}_{Q_k}q,
 \quad1\leq k\leq K,\qquad
 \mathfrak f=-\sum_{\alpha=1}^{g}v_\alpha s^{-n_\alpha}+O(s)
 \quad(P\to P_\infty).
\]
Both $\mathfrak f$ and the coefficients $v_\alpha$ are obtained by a
finite linear meromorphic interpolation problem. In genus zero the
central condition is $\mathfrak f=O(s)$ and the sum is empty.
\end{lemma}
\begin{proof}
For $g\geq1$, Riemann--Roch gives
$H^1(\mathcal O((2g-1)P_\infty))=0$. Thus arbitrary outer principal
parts extend to a meromorphic function with an additional central
pole of order at most $2g-1$ \cite[Chapter~IV]{Hartshorne1977}.
Subtract functions with their only pole at $P_\infty$ to remove the
nongap powers in descending order, then subtract the central constant.
A difference of two normalized solutions would have its only pole
at $P_\infty$ with largest order a gap, which is impossible; its zero
constant then makes it vanish. On the sphere the outer principal
parts extend without a central pole, and fixing the central value
gives the same conclusion.
\end{proof}

Use $1/\nu$ as coordinate at the outer points. Put $d_i=1$ for
$1\leq i\leq L$ and $d_j=-1$ for $L+1\leq j\leq K$, and define
\begin{equation}\label{eq:BA-higher-time-phase}
 p(s^{-1})=\sum_{\alpha=1}^{g}v_\alpha s^{-n_\alpha},\qquad
 b_k^{(r,\nu_0)}=-d_k(\mathfrak f+q)(Q_k),\qquad
 c_k(t)=c_k e^{-it b_k^{(r,\nu_0)}},\quad1\leq k\leq K.
\end{equation}
The sum $\mathfrak f+q$ is regular at every $Q_k$. Its values and
the coefficients of $p$ depend only on the fixed marked spectral data.
Let $\boldsymbol U_{n_\alpha}$ be the BA direction of the normalized
second-kind differential with principal part $d(s^{-n_\alpha})$.
The split into $p$ and mode phases uses the fixed local parameter $s$.
Changing that parameter changes this representation of the same
Hamiltonian evolution.

\begin{theorem}[The spectral BA hierarchy]\label{thm:BA-higher-flows}
Every prescribed realization in \cref{thm:general-BA-realization}
evolves under $H_{r,\nu_0}$ by the same ordinary and dual BA functions
and Cauchy kernel, with one-point time $p$ and gauges $c_k(t)$ from
\cref{eq:BA-higher-time-phase}. Explicitly, use those gauges in
\cref{eq:general-elementary-modes,eq:general-pair-modes}, and replace
the first-time exponentials by $e^{-itp}$ and $e^{itp}$ at the puncture.
The Abel argument is
\begin{equation}\label{eq:BA-higher-Abel-direction}
 z(t)=z_0-it\sum_{\alpha=1}^{g}v_\alpha\boldsymbol U_{n_\alpha}.
\end{equation}
The elementary partner phases are reciprocal; each pair mode has
phase $e^{it(b_i^{(r,\nu_0)}+b_j^{(r,\nu_0)})}$ and its partner
has the inverse phase. The statement holds locally on the nonspecial
BA chart, without a distinct-phase assumption.
Finite constant linear combinations of the Hamiltonians use the
corresponding sums of $p$ and the phase rates.

As $\nu_0$ varies, the flows with $2\leq r\leq K+1$ span all
Jacobian directions. Consequently finitely many commuting polynomial
spectral coefficients suffice to generate those translations.
\end{theorem}
\begin{proof}
Start with the prescribed BA modes and the right eigenvector
\cref{eq:BA-prescribed-eigenvector}. The left vector is
\[
 \ell(P)=\bigl((-c_i\Omega(P,Q_i))_{i=1}^{L},
 (\kappa_j\Omega(P,Q_j)/c_j)_{j=L+1}^{K},\psi(P)\bigr).
\]
The convolution identities give its spectral equation and
\begin{equation}\label{eq:BA-higher-initial-kernel}
 \ell(P)r(Q)=-(\nu(P)-\nu(Q))\Omega(P,Q).
\end{equation}
Evolve the matrix by the canonical Hamiltonian and these vectors by
$i\dot r_H=-\mathsf T_{r,\nu_0}r_H$ and
$i\dot\ell_H=\ell_H\mathsf T_{r,\nu_0}$. Polynomiality of the time
matrix preserves their finite pole divisors.

At infinity, write $\partial=i\partial_t$. The eigenvalue equation
gives the convergent local factorization
\begin{equation}\label{eq:BA-higher-local-factorization}
 r_H=e^{-itq(P)}\mathsf V(t,\nu(P))r(P,0),\qquad
 \partial\mathsf V=-\frac{\mathsf G}{\nu-\nu_0}\mathsf V,
 \qquad \mathsf V(0,\nu)=\Id.
\end{equation}
Since $\mathsf V=\Id+O(\nu^{-1})$, every leading coordinate direction
is preserved, even when branch intercepts coincide. Dress the vectors
reciprocally by
$\widetilde r=e^{-it\mathfrak f}r_H$ and
$\widetilde\ell=e^{it\mathfrak f}\ell_H$.
At the central point $q=O(\nu^{-1})$, so the last components have
the normalized dual and ordinary factors $e^{itp}$ and $e^{-itp}$;
all other components are $O(s)$ times those factors.
At each outer point the polar exponent cancels. If $R_k$ is the
finite part of $q+\mathfrak f$, the right vector's leading
coefficients are
\[
 -\nu e^{-itR_i}/c_i\quad(1\leq i\leq L),\qquad
 c_j\nu e^{-itR_j}/\kappa_j\quad(L+1\leq j\leq K).
\]
These are precisely $-\nu/c_i(t)$ and $c_j(t)\nu/\kappa_j$.

BA uniqueness identifies the last components. Each remaining component
has the fixed BA divisor, one additional simple pole at its own
outer point, the calculated residue and zero leading coefficient at
the central puncture. The same uniqueness identifies the kernel
components. No new poles arise at distinct points of equal $\nu$:
the derivative of $\ell_H(P)r_H(Q)$ is divisible by
$\nu(P)-\nu(Q)$, because the time matrix is polynomial. Thus all
initially removable poles in \cref{eq:BA-higher-initial-kernel}
remain removable. The diagonal pairing is constant, so the duality
differential and weights are unchanged. The central row and the
opposite outer blocks of the spectral equation now recover all modes
as in \cref{eq:BA-converse-recovery}. This proves their asserted
BA formulas, including every pair phase.

For the spanning assertion, choose an unramified finite fiber
$\nu=\nu_0$ with distinct eigenvalues $\lambda_h$, $1\leq h\leq K+1$.
The function $\mathfrak f+q$ has poles only at the central point,
with principal part $-p$, and on this fiber, with the principal parts
of $q$. For a holomorphic differential $du=w(P)d\nu$, the residue
theorem therefore gives
\[
 \operatorname{Res}_{P_\infty}(p\,du)
 =\sum_{h=1}^{K+1}\lambda_h^{r-1}w(P_h).
\]
If $du$ annihilates every such BA direction, the left side vanishes
for $2\leq r\leq K+1$. The same vanishing for $r=1$ follows by
applying the residue theorem to $du/(\nu-\nu_0)$.
The Vandermonde matrix is invertible, so $w$ vanishes on every such
fiber, hence $du=0$. Residue duality proves spanning.
Finite-dimensionality and expansion in $\nu_0$ give the final assertion.
\end{proof}

On a fixed spectral level the vectors in
\cref{eq:BA-higher-Abel-direction} are constant, and constant linear
combinations of the Hamiltonians select any Abel direction.
The required combination depends on the spectral data: the theorem
identifies polynomial Hamiltonian flows with level-dependent
combinations of normalized BA times.

\begin{corollary}[Higher flows from generic initial modes]
\label{cor:BA-higher-reconstruction}
On the generic initial-data domains of \cref{thm:BA-spectral-converse},
the same reconstructed curve and BA data solve the commuting flows
$H_{r,\nu_0}$ with \cref{eq:BA-higher-time-phase}.
For repeated phases, use the same constant active bases and restore
the original coordinates after the BA evolution.
\end{corollary}
\begin{proof}
Apply \cref{thm:BA-higher-flows} to the reconstructed data. Every
spectral Hamiltonian is invariant under equal-phase block conjugations,
so their moment matrices are conserved and the diagonalizing bases
remain constant. Inactive directions are direct summands of the
pencil; matrix powers preserve their zero couplings. Thus the reduction
and restoration in \cref{eq:BA-block-bases,eq:BA-inactive-splitting}
apply to all these flows.
\end{proof}

For the fourth trace, $q=\lambda^3/\nu$. Its finite part has a useful
modal-power interpretation. Write $e_k=\rho-b_k$ and let
$C_i=I_i-\sum_{j=L+1}^{K}I_{ij}$,
$C_j=I_j-\sum_{i=1}^{L}I_{ij}$, where
$I_k=\xi_k\bar\xi_k$ and $I_{ij}=\xi_{ij}\bar\xi_{ij}$.
On a prescribed realization, the eigenvector's pure coordinate pole
at each $Q_k$ makes the corresponding moment block diagonal in that
basis, also when phases coincide.
Second-order perturbation at infinity gives
\[
 \lambda=d_k\nu+e_k-d_kC_k/\nu+O(\nu^{-2}),\qquad
 q=d_k\nu^2+3e_k\nu+3d_k(e_k^2-C_k)+O(\nu^{-1}).
\]
Thus the phase in \cref{eq:BA-higher-time-phase} is computed from
the outer finite part of $\mathfrak f$ and conserved power charges.
On a fixed spectral level these constant phase rates can also be
removed by subtracting $\sum_{k=1}^{K}b_k^{(r,\nu_0)}C_k$ from the
Hamiltonian. The resulting flow has the normalized one-point BA time.

